%! Author = sriadityavedantam
%! Date = 3/30/26

\section{Homomorphisms and Isomorphisms}

\begin{definition}[Center]
    \[
        Z(G)\coloneqq \{x\in G: xa=ax\ \forall a\in G\}
        =
        \{x\in G: xax^{-1}=a\ \forall a\in G\}
        =
        \bigcap\limits_{a\in G} C_G(a).
    \]
\end{definition}

Note that the centralizer is only for one element $a$:
\[
    C_G(a)\coloneqq \{x\in G: xa=ax\}.
\]
Also,
\[
    Z(G)\trianglelefteq G.
\]
\begin{theorem}{
    Suppose $G/Z(G)$ is cyclic.
    Then $G$ is abelian.
    Hence $Z(G)=G$.
}
Denote $Z(G)=Z$.
Let $cZ$ be a generator for $G/Z$.
Then every coset in $G/Z$ has the form
\[
    aZ=c^{i}Z,
    \qquad
    bZ=c^{j}Z
\]
for some $i,j\in \z$.
Let $a\in c^{i}Z$ and $b\in c^{j}Z$.
Then there exist $d,e\in Z$ such that
\[
    a=c^{i}d,
    \qquad
    b=c^{j}e.
\]
Thus
\begin{align*}
    ab &= c^{i}dc^{j}e\\
    &= c^{i}c^{j}de\\
    &= c^{i}c^{j}ed\\
    &= c^{i+j}ed\\
    &= c^{j+i}de\\
    &= c^{j}c^{i}de\\
    &= c^{j}ec^{i}d\\
    &= ba.
\end{align*}
So $ab=ba$ for all $a,b\in G$.
Therefore $G$ is abelian.
Hence $Z(G)=G$.
\end{theorem}

\begin{definition}[Kernel of Group Functions]
    Suppose $\phi:G\to H$ is a group homomorphism.
    Define the kernel of $\phi$ by
    \[
        \ker\phi \coloneqq \{g\in G:\phi(g)=1_H\}.
    \]
\end{definition}

\begin{proposition}{
    \[
        \ker\phi\trianglelefteq G.
    \]
    Also, $\phi$ is injective if and only if
    \[
        \ker\phi=\{1_G\}.
    \]
}
    First, $1_G\in \ker\phi$ since
    \[
        \phi(1_G)=1_H.
    \]
    If $a,b\in \ker\phi$, then
    \[
        \phi(ab^{-1})=\phi(a)\phi(b)^{-1}=1_{H}1_H^{-1}=1_H.
    \]
    So $ab^{-1}\in \ker\phi$.
    Hence $\ker\phi\leq G$.
    Now let $g\in G$ and $x\in \ker\phi$.
    Then
    \[
        \phi(gxg^{-1})=\phi(g)\phi(x)\phi(g)^{-1}=\phi(g)1_H\phi(g)^{-1}=1_H.
    \]
    Thus
    \[
        gxg^{-1}\in \ker\phi.
    \]
    So
    \[
        \ker\phi\trianglelefteq G.
    \]
    Now suppose $\phi$ is injective.
    If $r\in \ker\phi$, then
    \[
        \phi(r)=1_H=\phi(1_G).
    \]
    Since $\phi$ is injective, $r=1_G$.
    Thus
    \[
        \ker\phi=\{1_G\}.
    \]
    Conversely, suppose
    \[
        \ker\phi=\{1_G\}.
    \]
    If $\phi(r)=\phi(s)$, then
    \[
        \phi(rs^{-1})=\phi(r)\phi(s)^{-1}=1_H.
    \]
    So
    \[
        rs^{-1}\in \ker\phi.
    \]
    Hence
    \[
        rs^{-1}=1_G,
    \]
    so $r=s$.
    Therefore $\phi$ is injective.
\end{proposition}

\begin{theorem}[First Isomorphism Theorem]{
    Suppose $\phi:G\to H$ is a group homomorphism.
    Let $K=\ker\phi$.
    Then
    \[
        G/K\cong \mathrm{Im}(\phi).
    \]
    Define
    \[
        \overline{\phi}:G/K\to \mathrm{Im}(\phi)
    \]
    by
    \[
        \overline{\phi}(gK)=\phi(g).
    \]
    Then $\overline{\phi}$ is a well-defined injective homomorphism whose image is $\mathrm{Im}(\phi)$.
}
To show $\overline{\phi}$ is well-defined, suppose
\[
    gK=hK.
\]
Then
\[
    h^{-1}g\in K=\ker\phi.
\]
So
\[
    \phi(h^{-1}g)=1_H.
\]
Hence
\[
    \phi(h)^{-1}\phi(g)=1_H,
\]
which implies
\[
    \phi(g)=\phi(h).
\]
Thus $\overline{\phi}(gK)=\overline{\phi}(hK)$.
Now for $gK,hK\in G/K$,
\[
    \overline{\phi}((gK)(hK))
    =
    \overline{\phi}(ghK)
    =
    \phi(gh)
    =
    \phi(g)\phi(h)
    =
    \overline{\phi}(gK)\overline{\phi}(hK).
\]
So $\overline{\phi}$ is a homomorphism.
It is injective because if
\[
    \overline{\phi}(gK)=1_H,
\]
then
\[
    \phi(g)=1_H,
\]
so $g\in K$.
Hence $gK=K$.
Thus
\[
    \ker(\overline{\phi})=\{K\},
\]
so $\overline{\phi}$ is injective.
Its image is exactly $\mathrm{Im}(\phi)$ by definition.
Therefore
\[
    G/K\cong \mathrm{Im}(\phi).
\]
\end{theorem}

Let
\[
    f:G/K\to \mathrm{Im}(\phi)
\]
be given by
\[
    f(gK)=\phi(g).
\]
This is the isomorphism from the theorem.

\begin{theorem}{
    The map
    \[
        g\mapsto \phi_g
    \]
    is a homomorphism from $G$ to $\Aut(G)$, where
    \[
        \phi_g(x)=gxg^{-1}.
    \]
}
Let $a,b\in G$.
We need to check that
\[
    \phi_{ab}=\phi_a\circ \phi_b.
\]
Let $x\in G$.
Then
\begin{align*}
    \phi_{ab}(x)
    &=abx(ab)^{-1}\\
    &=abxb^{-1}a^{-1}\\
    &=a\phi_b(x)a^{-1}\\
    &=\phi_a(\phi_b(x)).
\end{align*}
So
\[
    \phi_{ab}=\phi_a\circ \phi_b.
\]
Hence
\[
    g\mapsto \phi_g
\]
is a homomorphism.
Now compute its kernel.
We have
\[
    \ker\phi
    =
    \{y\in G:\phi_y(x)=x\ \forall x\in G\}.
\]
This means
\[
    yxy^{-1}=x\quad \forall x\in G.
\]
Equivalently,
\[
    yx=xy\quad \forall x\in G.
\]
So
\[
    \ker\phi=Z(G).
\]
\end{theorem}

\begin{corollary}{
    \[
        G/Z(G)\cong \Inn(G).
    \]
}
    This follows from the first isomorphism theorem applied to the homomorphism
    \[
        G\to \Aut(G),
        \qquad
        g\mapsto \phi_g.
    \]
    Its kernel is $Z(G)$, and its image is $\Inn(G)$.
    Hence
    \[
        G/Z(G)\cong \Inn(G).
    \]
\end{corollary}

If $G$ is abelian, then
\[
    Z(G)=G
\]
and
\[
    \Inn(G)=\{\id\}.
\]
If $G=S_n$ with $n\geq 3$, then
\[
    Z(G)=\{1\},
\]
so
\[
    G\cong \Inn(G)\leq \Aut(G).
\]
Thus
\[
    \Inn(G)\trianglelefteq \Aut(G).
\]
\begin{theorem}{
    Suppose
    \[
        K\trianglelefteq N\trianglelefteq G
    \]
    with also
    \[
        K\trianglelefteq G.
    \]
    Then
    \[
        N/K\trianglelefteq G/K
    \]
    and
    \[
        (G/K)/(N/K)\cong G/N.
    \]
}
Define
\[
    \phi:G/K\to G/N
\]
by
\[
    \phi(gK)=gN.
\]
First check that $\phi$ is well-defined.
Suppose
\[
    aK=bK.
\]
Then
\[
    b^{-1}a\in K.
\]
Since
\[
    K\leq N,
\]
we have
\[
    b^{-1}a\in N.
\]
Thus
\[
    aN=bN.
\]
Now compute the kernel.
Suppose
\[
    gK\in \ker\phi.
\]
Then
\[
    \phi(gK)=N.
\]
So
\[
    gN=N,
\]
which means
\[
    g\in N.
\]
Hence
\[
    gK\in N/K.
\]
Thus
\[
    \ker\phi=N/K.
\]
Also, $\phi$ is surjective, since for any $gN\in G/N$,
\[
    \phi(gK)=gN.
\]
Therefore, by the first isomorphism theorem,
\[
    (G/K)/(N/K)\cong G/N.
\]
Since $N/K=\ker\phi$, it is normal in $G/K$.
Thus
\[
    N/K\trianglelefteq G/K.
\]
\end{theorem}

Normality is not transitive.
It is possible that
\[
    K\trianglelefteq N
    \quad\text{and}\quad
    N\trianglelefteq G
\]
but
\[
    K\ntrianglelefteq G.
\]
For example, if
\[
    G=\z,\qquad N=6\z,\qquad K=12\z,
\]
then
\[
    K\trianglelefteq N\trianglelefteq G
    \quad\text{and}\quad
    K\trianglelefteq G.
\]
Also,
\begin{gather*}
    G/N=\z/6\z\cong \z_6,\\
    N/K=6\z/12\z\cong \z_2,\\
    G/K=\z/12\z\cong \z_{12},\\
\end{gather*}
and
\[
    (G/K)/(N/K)\cong \z_6\cong G/N.
\]

\begin{theorem}[4th Isomorphism Theorem]{
    Suppose
    \[
        N\trianglelefteq G.
    \]
    Then there is a bijection between subgroups of $G/N$ and subgroups of $G$ containing $N$.
}
Suppose $H\leq G$ with $N\leq H$.
Define
\[
    T=\{hN:h\in H\}.
\]
We check that $T\leq G/N$.
Since $1\in H$, we have
\[
    N=1N\in T.
\]
If $h_1,h_2\in H$, then
\[
    (h_{1}N)(h_{2}N)=(h_{1}h_2)N\in T.
\]
If $h\in H$, then
\[
    (hN)^{-1}=h^{-1}N\in T.
\]
So $T\leq G/N$.
Conversely, suppose $T\leq G/N$.
Define
\[
    H=\{h\in G:hN\in T\}.
\]
We check that $H\leq G$.
Since $N\in T$, we have $1\in H$.
If $h_1,h_2\in H$, then
\[
    h_{1}N,h_{2}N\in T.
\]
Since $T$ is a group,
\[
    (h_{1}N)(h_{2}N)=(h_{1}h_2)N\in T.
\]
So $h_{1}h_2\in H$.
If $h\in H$, then
\[
    hN\in T,
\]
so
\[
    (hN)^{-1}=h^{-1}N\in T.
\]
Thus $h^{-1}\in H$.
Hence $H\leq G$.
These maps are inverse to each other, so the correspondence is a bijection.
\end{theorem}

These maps are also inverses of one another.
If we start with $H\leq G$ containing $N$, then
\[
    T=\{hN:h\in H\}.
\]
If we start with $T\leq G/N$, then
\[
    H=\{h\in G:hN\in T\}.
\]

\begin{proposition}{
    Under these maps,
    \[
        H\trianglelefteq G
        \iff
        T\trianglelefteq G/N.
    \]
}
($\implies$).
    Let $xN\in T$ and $gN\in G/N$.
Then
    \[
        (gN)(xN)(g^{-1}N)=gxg^{-1}N.
    \]
    Since $H\trianglelefteq G$ and $x\in H$, we have
    \[
        gxg^{-1}\in H.
    \]
    Thus
    \[
        gxg^{-1}N\in T.
    \]
    So
    \[
        T\trianglelefteq G/N.
    \]

    ($\impliedby$).
    Suppose $T\trianglelefteq G/N$.
    Let $x\in H$ and $g\in G$.
    Since $xN\in T$, we have
    \[
        (gN)(xN)(g^{-1}N)\in T.
    \]
    That is,
    \[
        gxg^{-1}N\in T.
    \]
    Hence
    \[
        gxg^{-1}\in H.
    \]
    So
    \[
        H\trianglelefteq G.
    \]
\end{proposition}

\begin{definition}[Simple Group]
    Suppose $G$ is a group and
    \[
        \{1\}\neq G.
    \]
    We say $G$ is simple if and only if
    \[
        N\trianglelefteq G
    \]
    implies
    \[
        N=G
        \quad\text{or}\quad
        N=\{1\}.
    \]
\end{definition}

\begin{theorem}{
    If $G$ is abelian and simple, then
    \[
        G\cong C_p
    \]
    for some prime $p$.
}
Since $G$ is abelian, every subgroup is normal.
Because $G$ is simple, it has no proper nontrivial subgroups.
Take
\[
    1\neq a\in G.
\]
Then
\[
    \gen{a}\leq G.
\]
Since $\gen{a}\neq \{1\}$ and $G$ is simple, we must have
\[
    \gen{a}=G.
\]
Thus $G$ is cyclic.So
\[
    G\cong C_n
\]
for some $n$.
If $n$ were composite, then there would exist a proper nontrivial subgroup of $C_n$, namely
\[
    \gen{a^k}
\]
for a suitable divisor $k$ of $n$.
But that would contradict simplicity.
Therefore $n$ must be prime.
Hence
\[
    G\cong C_p
\]
for some prime $p$.
\end{theorem}

What does it mean when a group is not simple.
It means it has a proper normal subgroup.
Suppose $G$ is not simple.
Then
\[
    1\neq N\triangleleft G.
\]
Thus one studies the quotient $G/N$ and can lift proper normal subgroups of $G/N$ back to $G$.
Now the correct thing to ask is what lift means.

\begin{definition}[Lift]
    By the 4th Isomorphism Theorem, there is a bijection between subgroups
    \[
        T\leq G/N
    \]
    and subgroups
    \[
        H\leq G
    \]
    containing $N$.
    We lift $T$ by assigning to it its corresponding subgroup $H$ in $G$.
\end{definition}

\begin{definition}[Jordan--Holder Series / Composition Series]
    If $G$ is not simple, we may find a proper normal subgroup $N\triangleleft G$.
    Then we may continue by finding a proper normal subgroup of $G/N$, and so on.
    Eventually we run out of proper normal subgroups.
    This process creates a sequence
    \[
        \{1\}\leq G_1\leq G_2\leq \dots \leq G_n\leq G
    \]
    with
    \[
        G_i\trianglelefteq G_{i+1}
    \]
    and each quotient
    \[
        G_{i+1}/G_i
    \]
    simple.
    This sequence is called a Jordan--Holder series, or composition series.
\end{definition}

\begin{example}
    Let
    \[
        G=C_6\cong C_3\times C_2.
    \]
    Then
    \[
        N_1=C_3\times \{1\}\trianglelefteq G
    \]
    and
    \[
        N_2=\{1\}\times C_2\trianglelefteq G.
    \]
    Thus
    \[
        G/N_1\cong C_2
    \]
    and
    \[
        G/N_2\cong C_3.
    \]
    So we have two composition series:
    \[
        \{1\}\trianglelefteq N_1\trianglelefteq G
    \]
    and
    \[
        \{1\}\trianglelefteq N_2\trianglelefteq G.
    \]
    Their composition factors are the same up to order:
    \[
        N_1/\{1\}\cong C_3,\qquad G/N_1\cong C_2,
    \]
    and
    \[
        N_2/\{1\}\cong C_2,\qquad G/N_2\cong C_3.
    \]
\end{example}